\documentclass[12pt,a4paper]{article}
\usepackage[utf8]{inputenc}
\usepackage[spanish]{babel}
\usepackage{amsmath}
\usepackage{amsfonts}
\usepackage{amssymb}
\usepackage{graphicx}
\usepackage[left=2cm,right=2cm,top=2cm,bottom=2cm]{geometry}
\usepackage{float}
\usepackage{hyperref}
\usepackage{xcolor}
\usepackage{titlesec}
\usepackage{subcaption}

\definecolor{twitchpurple}{RGB}{145,70,255}
\definecolor{darkgray}{RGB}{64,64,64}

\titleformat{\section}{\Large\bfseries\color{twitchpurple}}{\thesection}{1em}{}
\titleformat{\subsection}{\large\bfseries\color{darkgray}}{\thesubsection}{1em}{}

\begin{document}

\begin{titlepage}
    \centering
    % Imagen de Twitch centrada al inicio
    
\includegraphics[width=0.25\textwidth]{./data/Twitch.png} \\
    \vspace{1cm}
    
    {\Huge\bfseries\color{twitchpurple} Análisis Multidimensional del Ecosistema Twitch} \\
    \vspace{0.5cm}
    {\Large Estudio Integral de Contenido, Creadores y Rendimiento Económico} \\
    \vspace{0.3cm}
    {\large Período 2016--2025}
    
    \vfill
    
    % Autores en la parte inferior
    \textbf{Grupo 4 - Visualización de Datos} \\
    \vspace{0.2cm}
    Alex Urrutia \quad $\bullet$ \quad Ariel Pulgar \quad $\bullet$ \quad Dante Aspee \\
    \vspace{0.1cm}
    \textit{USM - Visualización de Datos} \\
    \vspace{0.1cm}
    \href{https://github.com/AlexUrrutiaUsm/VisualisacionGrupo4SJ}{Repositorio GitHub}\\
    \vspace{0.1cm}
    \today
\end{titlepage}

\newpage
\tableofcontents
\newpage

\section{Introducción} 

\subsection{Contexto y Relevancia} 

Twitch se ha consolidado como la plataforma de streaming dominante en el ecosistema gaming y entretenimiento digital. Su crecimiento exponencial desde 2016 ha transformado no solo la industria de los videojuegos, sino también las estrategias de marketing, el comportamiento del consumidor digital y las nuevas economías creativas.

\subsection{Objetivos del Análisis}

El presente estudio adopta un enfoque multidimensional para comprender: 

\begin{itemize}
    \item \textbf{Dimensión de Contenido}: Evolución de las preferencias de juegos y patrones de consumo 
    \item \textbf{Dimensión de Creadores}: Dinámicas de crecimiento y comportamiento de streamers exitosos 
    \item \textbf{Dimensión Económica}: Impacto financiero y métricas de audiencia de la plataforma 
\end{itemize}

\subsection{Metodología Integrada} 

Para garantizar la coherencia analítica, cada dimensión del estudio utiliza: 
\begin{enumerate}
    \item Datos que abarcan el período 2016--2025 
    \item Métricas estandarizadas de \textit{engagement} y rendimiento 
    \item Enfoque en los elementos más representativos de cada categoría
\end{enumerate}
\newpage
\section{Análisis de la Evolución del Contenido Gaming} 
\textit{Sección realizada por: Alex Urrutia}

\subsection{Marco Conceptual: Dinámica Stream vs Watch} 

El primer eje de análisis examina la relación entre la popularidad de streaming (cantidad de horas transmitidas) y el consumo real (horas visualizadas) de los principales videojuegos en Twitch. Esta perspectiva revela discrepancias importantes entre la oferta de contenido y la demanda real de audiencia.

\subsection{Metodología: Slope Chart Analysis}

El análisis utiliza un slope chart comparativo que visualiza simultáneamente: 
\begin{itemize}
    \item \textbf{Stream Rank}: Posicionamiento por horas totales transmitidas 
    \item \textbf{Watch Rank}: Posicionamiento por horas totales visualizadas 
    \item \textbf{Trayectorias de cambio}: Movimiento relativo entre ambas métricas 
\end{itemize}

\begin{figure}[H]
    \centering
    \includegraphics[width=0.9\textwidth]{"./Alex_Urrutia/Grafico Streamed Vs Watched.png"}
    \caption{Slope Chart: Stream Rank vs Watch Rank - Top Juegos en Twitch (2016--2025)}
\label{fig:slope_chart}
\end{figure}

\subsection{Hallazgos Principales} 

\begin{itemize}
    \item \textbf{Juegos que mantienen equilibrio}: Títulos como \textit{Just Chatting} y \textit{League of Legends} mantienen rankings similares en ambas dimensiones, indicando equilibrio entre oferta y demanda.
    
    \item \textbf{Juegos con más streamers que espectadores}: Juegos como \textit{Rainbow Six Siege} y \textit{Apex Legends} muestran mayor ranking de stream que de watch, sugiriendo saturación de contenido.
    
    \item \textbf{Juego \textit{outliner}}: \textit{Hearthstone} muestra un ranking de watch significativamente superior al de stream, indicando alta demanda con oferta limitada.
\end{itemize}

\section{Análisis del Ecosistema de Creadores de Contenido}
\textit{Sección realizada por: Ariel Pulgar}

\subsection{Marco Conceptual: Jerarquías de Influencia} 

El segundo eje examina la distribución y características de los streamers más exitosos, enfocándose en dos métricas fundamentales: base de seguidores establecida y capacidad de crecimiento activo. Este análisis revela las dinámicas de poder e influencia dentro del ecosistema Twitch, donde se utilizó un gráfico de nube de palabras, ya que es un gráfico llamativo el cual se puede modificar la ubicación de las palabras de forma estratégica, por ejemplo, los datos que utilice son de una plataforma de stream llamada “Twitch”, las palabras se podrían ordenar de distintas orientaciones y ubicaciones para formar el logo de la plataforma (no se realizó este proceso, pero creo que sería una idea muy llamativa).
\subsection{Metodología: Word Cloud Ponderado} 

La visualización utiliza una nube de palabras donde: 
\begin{itemize}
    \item \textbf{Tamaño de texto}: Proporcional al número total de seguidores 
    \item \textbf{Intensidad de color}: Representa la tasa de seguidores ganados (crecimiento) 
    \item \textbf{Posicionamiento}: Refleja la combinación de ambas métricas 
\end{itemize}

\begin{figure}[H]
    \centering
    \includegraphics[width=0.9\textwidth]{"./Ariel_Pulgar/nube_palabras_followers.png"}
    \caption{Nube de Palabras: Streamers por Seguidores y Crecimiento}
\label{fig:wordcloud}
\end{figure}

\newpage

\section{Análisis del Impacto Económico y Audiencia}
\textit{Sección realizada por: Dante Aspee}

\subsection{Metodología y Justificación de Criterios}

Para el desarrollo de esta dimensión, he seleccionado como criterios de análisis los \textbf{Ingresos Anuales Totales} y la \textbf{Audiencia Promedio}. Se han seleccionado estos criterios con el fin de obtener una lectura objetiva sobre la viabilidad económica y el alcance real de la plataforma en el mercado global de entretenimiento digital:

\begin{itemize}
    \item \textbf{Ingresos Anuales}: Este criterio permite cuantificar el éxito de Twitch como modelo de negocio y su capacidad para capturar valor económico dentro de la industria del entretenimiento digital. Es la métrica definitiva para evaluar la sostenibilidad financiera a largo plazo.
    \item \textbf{Audiencia Promedio}: Se seleccionó este criterio porque representa la materia prima de la plataforma. Sin una base de usuarios sólida, el valor de los creadores y el interés de los anunciantes disminuye. Por lo tanto es fundamental entender no solo cuánto dinero se genera, sino cuántas personas están consumiendo el contenido que Twitch ofrece.
\end{itemize}

La relevancia de analizar ambos criterios en conjunto radica en que permite identificar la \textbf{eficiencia de monetización}. No basta con saber cuánto dinero ingresa, sino entender si dicho ingreso es proporcional al volumen de personas consumiendo contenido o si la plataforma está logrando extraer más valor de una audiencia más pequeña a través de suscripciones y publicidad diversificada.

En cuanto a la grafica seleccionada, la visualización emplea un \textbf{Lollipop Chart Dual} que representa:
\begin{itemize}
    \item \textbf{Eje Y}: Ingresos en miles de millones de USD.
    \item \textbf{Hexágonos}: Audiencia promedio (millones de espectadores).
    \item \textbf{Línea temporal}: Evolución del periodo 2016--2024.
\end{itemize}

\begin{figure}[H]
    \centering
    \includegraphics[width=0.9\textwidth]{"./Dante_Aspee/Lollipop_Twitch.png"}
    \caption{Desempeño Financiero y Audiencia de Twitch (2016--2024). Fuente: Elaboración propia a partir de datos de Twitch Tracker, Business of Apps y Demand Sage \cite{DemandSage}.}
    \label{fig:lollipop}
\end{figure}

\subsection{Análisis de Resultados y Conclusiones}

Tras el análisis de los datos presentados en la Figura \ref{fig:lollipop}, se desprenden las siguientes conclusiones clave:

\begin{itemize}
    \item \textbf{Fase de Expansión (2016--2021)}: Se identifica una correlación directa entre el aumento masivo de espectadores y el crecimiento de ingresos, fenómeno potenciado significativamente por la pandemia de COVID-19 y el auge de los eventos competitivos globales de eSport.
    \item \textbf{Fase de Maduración (2021--2024)}: Se observa un fenómeno de estabilización. A pesar de una leve disminución en la audiencia promedio post-pandemia (de 2.8M a 2.4M), los ingresos se han mantenido resilientes por encima de los \$1.8B USD. Esto confirma que la plataforma ha transitado exitosamente de un crecimiento basado en volumen a uno basado en la optimización económica de su base de usuarios existente, a través de estrategias de monetización diversificadas como suscripciones, publicidad y acuerdos de patrocinio.
\end{itemize}
\newpage
\section{Conexión entre los puntos analizados}
\subsection{¿Cómo se relacionan las partes del estudio?} 

Este análisis muestra que el éxito de Twitch no ocurre por casualidad, sino por cómo se conectan los tres niveles estudiados contenido, creadores y economía. Cada uno de estos elementos influye en los otros, creando un ciclo de retroalimentación que impulsa el crecimiento y la sostenibilidad de la plataforma.

\subsubsection{Relación entre juegos y streamers} 
Elegir el juego correcto es clave para un creador. El estudio indica que los juegos con un mejor equilibrio entre quienes transmiten y quienes ven el contenido ayudan a que los streamers tengan un crecimiento más constante. No se trata solo de transmitir, sino de encontrar dónde hay demanda real de audiencia.

\subsubsection{El peso de los grandes creadores} 
Los streamers con más seguidores y crecimiento activo son los que generan la mayor parte del movimiento en la plataforma. Aunque son un grupo específico, su capacidad para atraer público concentra gran parte del valor total de Twitch.

\subsubsection{El resultado económico final} 
El crecimiento del dinero que gana Twitch depende directamente de que existan tanto juegos variados como creadores que impacten a la audiencia. Si estos dos elementos funcionan, la plataforma logra resultados económicos sólidos.

\subsection{La clave del éxito: Una visión integrada} 

Podemos entender el éxito de Twitch como el resultado de tres factores que deben trabajar en conjunto. Primero, es fundamental la \textbf{eficiencia del contenido}, que ocurre cuando se eligen categorías de juegos con un ratio ideal entre oferta y demanda. A esto se suma el \textbf{crecimiento de los creadores}, que depende de saber equilibrar la base de seguidores que ya tienen con la llegada de gente nueva. Finalmente, todo esto permite una \textbf{optimización económica}, que es la capacidad de transformar ese interés de la gente en ingresos reales de forma eficiente. En definitiva, la plataforma funciona mejor cuando logra que estas tres piezas encajen al mismo tiempo. Por eso, el análisis multidimensional es tan importante: nos permite ver cómo cada parte influye en las otras y entender que el éxito de Twitch no es solo cuestión de tener buenos juegos o grandes streamers, sino de cómo todos estos elementos se conectan para crear un ecosistema sólido y rentable.

\newpage
\section{Gráfico generado por IA Generativa}

Para explorar el potencial del uso de inteligencia artificial generativa en el análisis de datos, se le entregó a \textbf{Claude} (Anthropic) acceso completo al repositorio, solicitándole que generara una visualización original. El prompt entregado fue el siguiente:

\begin{quote}
\textit{``Por favor, analisa el repositorio de forma completa, para luego, crees un archivo .py en la carpeta Claude con la toda la informacion de la carpeta data quiero que dicha visualisacion sea con respecto a un criterio que encuentres relevante analizar, ademas has que dicha visualisacion sea con un grafico poco conocido y diferente a los graficos ya utilizados por favor.''}
\end{quote}

Como resultado, el modelo generó un \textbf{Stream Graph} (gráfico de flujo apilado), una visualización poco convencional que permite observar la evolución temporal de múltiples categorías de forma fluida y orgánica. El criterio elegido por la IA fue el análisis de las \textbf{horas vistas por los juegos más populares en Twitch} a lo largo del período 2016--2025, destacando tendencias de crecimiento, dominancia y declive de cada título. Luego de un ajuste menor en la leyenda, que se salía levemente de su zona delimitada, el gráfico final obtenido fue el siguiente:

\begin{figure}[H]
    \centering
    \includegraphics[width=0.95\textwidth]{"./Claude/streamgraph_twitch.png"}
    \caption{Stream Graph generado por Claude: Evolución de horas vistas por juego en Twitch (2016--2025).}
    \label{fig:streamgraph}
\end{figure}

\newpage
\section{Conclusiones y Perspectivas Futuras}

\subsection{Síntesis de Hallazgos}

En síntesis, el análisis multidimensional del ecosistema Twitch entre 2016 y 2025 demuestra que el éxito de la plataforma depende de la interacción constante de tres dimensiones inseparables:

\begin{itemize}
    \item \textbf{Contenido}: Se identificó un catálogo donde títulos como \textit{Just Chatting} y \textit{League of Legends} logran equilibrar la oferta y la demanda. En contraste, otros juegos presentan saturación o una demanda insatisfecha que representa oportunidades de crecimiento.
    \item \textbf{Creadores}: Existe una capa altamente concentrada donde un grupo reducido de streamers, caracterizados por una gran base de seguidores y un crecimiento activo, concentra la mayor parte del valor y la audiencia de la plataforma.
    \item \textbf{Economía}: La dimensión financiera transitó de una fase de expansión acelerada por la pandemia y los eSports hacia una etapa de maduración. En esta fase, la monetización eficiente sostiene ingresos superiores a los \$1.8B USD, incluso ante una audiencia que se ha estabilizado levemente tras el pico de 2021.
\end{itemize}

Este estudio se complementó con un experimento de IA generativa, el cual demostró que herramientas como Claude pueden identificar criterios de análisis relevantes de forma autónoma. La creación de visualizaciones no convencionales, como el \textit{Stream Graph}, enriqueció la comprensión del ecosistema desde ángulos que el análisis humano no siempre prioriza.

\subsection{Visión y Perspectivas Finales}

Este trabajo ofrece un modelo para entender cómo el contenido y los creadores impulsan la economía digital. Como pasos siguientes, se propone explorar el análisis predictivo de nuevas tendencias de contenido y realizar estudios comparativos con otras plataformas de streaming para entender mejor la competencia global en la industria.

\newpage

\section*{Referencias} 

\begin{itemize}
    \item \textbf{SullyGnome}. \textit{Twitch Statistics and Analytics}. Recuperado de \url{https://sullygnome.com}\\
    \item \textbf{TwitchTracker}. \textit{Twitch Channel Statistics and Charts}. Recuperado de \url{https://twitchtracker.com}\\
    \item \textbf{Business of Apps}. \textit{Twitch Revenue and Usage Statistics}. Recuperado de \url{https://www.businessofapps.com/data/twitch-statistics/}\\
    \item \textbf{Demand Sage}. \textit{Twitch User Count and Demographics}. Recuperado de \url{https://www.demandsage.com/twitch-users/}
\end{itemize}

\vspace{0.5cm}
\small \textbf{Nota:} Debido a que Amazon no desglosa oficialmente los ingresos de Twitch, los datos financieros representan estimaciones consolidadas de la industria validadas por las fuentes citadas.

\end{document}